Finally, rather than incorporate structural coherence into the model, we treat it as an
evaluation measure that identifies which of the recovered ideologies relate more strongly
to their own concepts than to concepts outside them. For ideology $k$ with support
$S_k=\{u:\max_f\eta_{k,f,u}>0\}$, we compare the association weight among concepts inside
$S_k$ to the weight crossing its boundary:
\[
  C_k \;=\;
  \frac{\sum_{(u,u')\in\Pcal,\; u\in S_k,\; u'\in S_k} W_{k,uu'}}
       {\sum_{(u,u')\in\Pcal,\; u\in S_k \,\oplus\, u'\in S_k} W_{k,uu'}}.
\]
Up to reparameterization this is the weighted conductance of $S_k$
\citep{shi2000,kannan2004}, since $\phi(S_k)=c_{S_k}/(2m_{S_k}+c_{S_k})=1/(1+2C_k)$ for
internal weight $m_{S_k}$ and cut $c_{S_k}$; we report the ratio rather than $\phi$ only
because $C_k=1$ marks the point the framework's definition asserts, namely as much
association weight inside an ideology as crossing its boundary. We use a different null
than is standard for such measures: instead of a configuration model, we compare $C_k$
against $200$ random supports matched to $S_k$ in size and in composition, sampled
stratified by concept type and, within type, by weighted-degree quartile in $W_k$.
Composition matters here because $\CC5$ leaves some rung pairs structurally incapable of
carrying weight, so a support drawn from a single rung would score well whatever concepts
it contained. Discarding draws for which the ratio is undefined, we report
$p=\#\{C_k^{\text{null}}\ge C_k\}/\#\{\text{finite }C_k^{\text{null}}\}$, whose smallest
attainable value at $200$ permutations is $.005$; further detail is in
Appendix~\ref{app:coherence}.
